%! Graphing Polynomial Functions and Modeling — Unit 4, Lesson 4.6 — Group Activity (Answer Key)
\documentclass[10pt]{article}
\usepackage{algebra2-article}
\usepackage{algebra2-key}   % pulls in -boxes; do NOT also load -boxes
\newcommand{\TallMath}[1]{$\displaystyle #1\rule[-1.4em]{0pt}{3.2em}$}

\newcommand{\lgrid}[4]{%
  \draw[step=1, linegray!40, very thin] (#1,#3) grid (#2,#4);
  \draw[->, semithick, charcoal] (#1,0)--(#2,0) node[below right, font=\scriptsize]{$x$};
  \draw[->, semithick, charcoal] (0,#3)--(0,#4) node[left, font=\scriptsize]{$y$};}

\begin{document}

\pageheader{Unit 4, Lesson 4.6}{Group Activity --- Answer Key}
\namepartnerperiod

\vspace{0.08in}

\begin{headlinebox}{forestbg}
\small\textbf{Predict, Then Look.} Every problem on this page runs the same five steps: \textbf{(1)}
degree and leading coefficient give the arms, \textbf{(2)} the factors give the $x$-intercepts,
\textbf{(3)} the multiplicities say cross, touch, or flatten, \textbf{(4)} $f(0)$ gives one exact
point, \textbf{(5)} one test value per interval says above or below. Make the prediction
\emph{before} you look at any graph --- the point of the page is that the equation already told you.
\end{headlinebox}

\vspace{0.06in}

\begin{tcolorbox}[colback=white, colframe=black!40, title=\textbf{Tier R --- Remediate},
    fonttitle=\bfseries, arc=1.5mm, left=3mm, right=3mm, top=2mm, bottom=2mm, breakable]
\begin{enumerate}[leftmargin=1.7em, itemsep=9pt]
  \item \textbf{Three cousins.} These three functions use the same two numbers, $-1$ and $2$. Fill in
    the table from the factored form --- no expanding.

    \vspace{3pt}
    \renewcommand{\arraystretch}{1.5}
    \begin{tabularx}{\linewidth}{|X|C{1.2cm}|C{1.4cm}|C{2.6cm}|C{2.6cm}|C{1.6cm}|}
    \hline
    \rowcolor{forestbg}
    \textbf{Function} & \textbf{Deg.} & \textbf{Lead} & \textbf{Left arm / right arm} & \textbf{At $x=-1$ / at $x=2$} & \textbf{$y$-int} \\ \hline
    (a) $f(x)=(x+1)(x-2)^2$  & \ans{$3$} & \ans{$1$}  & \ans{down / up} & \ans{cross / touch} & \ans{$4$} \\ \hline
    (b) $g(x)=-(x+1)(x-2)^2$ & \ans{$3$} & \ans{$-1$} & \ans{up / down} & \ans{cross / touch} & \ans{$-4$} \\ \hline
    (c) $h(x)=(x+1)^2(x-2)$  & \ans{$3$} & \ans{$1$}  & \ans{down / up} & \ans{touch / cross} & \ans{$-2$} \\ \hline
    \end{tabularx}

  \item \textbf{Now match.} Each graph below is one of the three. Write the letter under each one.
    (Each horizontal gridline is $2$ units.)
    \begin{center}
    \begin{tikzpicture}[scale=0.42]
      \lgrid{-3}{3}{-4}{4}
      \begin{scope}
        \clip (-3,-4) rectangle (3,4);
        \draw[very thick, forest, domain=-1.6:2.7, samples=150, smooth]
          plot (\x,{(0 - (\x)*(\x)*(\x) + 3*(\x)*(\x) - 4)/2});
      \end{scope}
      \foreach \x in {-2,-1,1,2}
        \draw[charcoal] (\x,-0.15)--(\x,0.15) node[below=1pt, font=\tiny]{$\x$};
      \node[font=\scriptsize\bfseries, charcoal] at (-2.2,3.4) {I};
    \end{tikzpicture}
    \hspace{0.15cm}
    \begin{tikzpicture}[scale=0.42]
      \lgrid{-3}{3}{-4}{4}
      \begin{scope}
        \clip (-3,-4) rectangle (3,4);
        \draw[very thick, forest, domain=-1.6:2.7, samples=150, smooth]
          plot (\x,{((\x)*(\x)*(\x) - 3*(\x)*(\x) + 4)/2});
      \end{scope}
      \foreach \x in {-2,-1,1,2}
        \draw[charcoal] (\x,-0.15)--(\x,0.15) node[below=1pt, font=\tiny]{$\x$};
      \node[font=\scriptsize\bfseries, charcoal] at (-2.2,3.4) {II};
    \end{tikzpicture}
    \hspace{0.15cm}
    \begin{tikzpicture}[scale=0.42]
      \lgrid{-3}{3}{-4}{4}
      \begin{scope}
        \clip (-3,-4) rectangle (3,4);
        \draw[very thick, forest, domain=-2.2:2.3, samples=150, smooth]
          plot (\x,{((\x)*(\x)*(\x) - 3*(\x) - 2)/2});
      \end{scope}
      \foreach \x in {-2,-1,1,2}
        \draw[charcoal] (\x,-0.15)--(\x,0.15) node[below=1pt, font=\tiny]{$\x$};
      \node[font=\scriptsize\bfseries, charcoal] at (-2.2,3.4) {III};
    \end{tikzpicture}
    \end{center}

    \vspace{2pt}
    Graph I is \ans{(b)} \qquad Graph II is \ans{(a)} \qquad Graph III is \ans{(c)}

    \vspace{2pt}
    Which \emph{one} feature separates (a) from (c)? \ans{which zero is the double one --- (a) touches
    at $2$, (c) touches at $-1$}

  \item \textbf{Sign chart practice.} Use $f(x) = (x+1)(x-2)^2$ from row (a).

    \vspace{3pt}
    \renewcommand{\arraystretch}{1.5}
    \begin{tabularx}{\linewidth}{|C{2.6cm}|C{1.6cm}|C{3.6cm}|C{1.8cm}|X|}
    \hline
    \rowcolor{forestbg}
    \textbf{Interval} & \textbf{Test $x$} & \textbf{$f(\text{test }x) = $} & \textbf{Sign} & \textbf{Above / below} \\ \hline
    $x<-1$    & $-2$ & \ans{$(-1)(16)=-16$} & \ans{$-$} & \ans{below} \\ \hline
    $-1<x<2$  & $0$  & \ans{$(1)(4)=4$}     & \ans{$+$} & \ans{above} \\ \hline
    $x>2$     & $3$  & \ans{$(4)(1)=4$}     & \ans{$+$} & \ans{above} \\ \hline
    \end{tabularx}

    \vspace{2pt}
    The sign changed at $x=-1$ but \emph{not} at $x=2$. Why not? \ans{the factor $(x-2)$ is squared,
    so it never turns negative --- the graph touches and turns back}

  \item \textbf{Count the turns.} A polynomial has degree $5$.

    \vspace{2pt}
    At most how many turning points can its graph have? \ans{$4$} \quad At most how many
    $x$-intercepts? \ans{$5$}

    \vspace{2pt}
    Must its graph cross the $x$-axis at least once? \ans{yes} \; Why? \ans{odd degree sends the arms
    in opposite directions, so it has to pass through}
\end{enumerate}
\end{tcolorbox}

\begin{tcolorbox}[colback=white, colframe=black!40, title=\textbf{Tier A --- Approaching Proficiency},
    fonttitle=\bfseries, arc=1.5mm, left=3mm, right=3mm, top=2mm, bottom=2mm, breakable]
\begin{enumerate}[leftmargin=1.7em, itemsep=10pt]
  \item \textbf{Factor first, then predict.} $f(x) = x^3 - x^2 - 6x$

    \vspace{2pt}
    (a) Factor completely (GCF first --- Lesson 4.2): \ans{$x(x-3)(x+2)$}

    \vspace{2pt}
    (b) Degree \ans{$3$}, leading coefficient \ans{$1$}. End behavior: as $x\to-\infty$,
    $f(x)\to$ \ans{$-\infty$}; as $x\to+\infty$, $f(x)\to$ \ans{$+\infty$}

    \vspace{2pt}
    (c) Zeros: \ans{$-2,\,0,\,3$} \quad Each one is a \ans{crossing} (multiplicity $1$).
    \quad $y$-intercept: \ans{$(0,0)$}

    \vspace{2pt}
    (d) Sign chart --- four intervals this time:

    \vspace{2pt}
    \renewcommand{\arraystretch}{1.5}
    \begin{tabularx}{\linewidth}{|C{2.4cm}|C{1.4cm}|C{3.0cm}|C{1.6cm}|X|}
    \hline
    \rowcolor{forestbg}
    \textbf{Interval} & \textbf{Test $x$} & \textbf{$f(\text{test }x)$} & \textbf{Sign} & \textbf{Above / below} \\ \hline
    $x<-2$    & $-3$ & \ans{$-18$} & \ans{$-$} & \ans{below} \\ \hline
    $-2<x<0$  & $-1$ & \ans{$4$}   & \ans{$+$} & \ans{above} \\ \hline
    $0<x<3$   & $1$  & \ans{$-6$}  & \ans{$-$} & \ans{below} \\ \hline
    $x>3$     & $4$  & \ans{$24$}  & \ans{$+$} & \ans{above} \\ \hline
    \end{tabularx}

    \vspace{3pt}
    (e) Here is the real graph. (Each horizontal gridline is $4$ units.) Check every prediction you
    just made against it.
    \begin{center}
    \begin{tikzpicture}[scale=0.5]
      \lgrid{-3}{4}{-3}{3}
      \begin{scope}
        \clip (-3,-3) rectangle (4,3);
        \draw[very thick, forest, domain=-2.5:3.5, samples=170, smooth]
          plot (\x,{((\x)*(\x)*(\x) - (\x)*(\x) - 6*(\x))/4});
      \end{scope}
      \foreach \x in {-2,-1,1,2,3}
        \draw[charcoal] (\x,-0.15)--(\x,0.15) node[below=1pt, font=\tiny]{$\x$};
      \foreach \p in {-2,0,3} \fill[charcoal] (\p,0) circle (2.5pt);
    \end{tikzpicture}
    \end{center}

    \vspace{2pt}
    (f) From the graph: \textbf{domain} \ans{all reals}, \textbf{range} \ans{all reals}. \;
    Relative maximum near $x\approx$ \ans{$-1.1$}, relative minimum near $x\approx$ \ans{$1.8$}.

    \vspace{2pt}
    (g) Does this function have an \emph{absolute} maximum or minimum? \ans{no} \; Explain using
    the degree: \ans{odd degree --- one arm runs to $+\infty$ and the other to $-\infty$}

  \item \textbf{Read the equation off the picture.} The graph below is a cubic.
    (Each horizontal gridline is $4$ units.)
    \begin{center}
    \begin{tikzpicture}[scale=0.5]
      \lgrid{-4}{3}{-3}{3}
      \begin{scope}
        \clip (-4,-3) rectangle (3,3);
        \draw[very thick, forest, domain=-3.6:2.0, samples=170, smooth]
          plot (\x,{(0 - (\x)*(\x)*(\x) - (\x)*(\x) + 5*(\x) - 3)/4});
      \end{scope}
      \foreach \x in {-3,-2,-1,1,2}
        \draw[charcoal] (\x,-0.15)--(\x,0.15) node[below=1pt, font=\tiny]{$\x$};
      \fill[charcoal] (-3,0) circle (2.5pt);
      \fill[charcoal] (1,0)  circle (2.5pt);
    \end{tikzpicture}
    \end{center}

    \vspace{2pt}
    (a) $x$-intercepts: \ans{$-3$ and $1$} \; Which one is a touch? \ans{$x=1$} \; So which factor is
    squared? \ans{$(x-1)$}

    \vspace{2pt}
    (b) Skeleton: $f(x) = a\,($\ans{$x+3$}$)($\ans{$x-1$}$)^2$

    \vspace{2pt}
    (c) The graph passes through $(0,-3)$, so $a = $ \ans{$-1$}. \; (Check the sign against the
    arms: is the left arm up or down? \ans{up})

    \vspace{2pt}
    (d) Final equation: \ans{$f(x) = -(x+3)(x-1)^2$}

  \item \textbf{Error analysis.} Two classmates, two different mistakes.

    \vspace{2pt}
    (a) ``The graph of $y=(x-4)^2(x+1)$ crosses the $x$-axis at $x=4$ and at $x=-1$.''
    What is wrong, and what does the graph actually do at $x=4$? \ans{$x=4$ has multiplicity $2$, so
    the graph touches there and turns back; only $x=-1$ is a crossing}

    \vspace{2pt}
    (b) ``This graph has four turning points, so the function has degree $4$.''
    What is the smallest degree it could have? \ans{$5$} \; Explain: \ans{degree $n$ allows at most
    $n-1$ turns, so $4$ turns needs $n\ge5$ --- and the graph only gives a \emph{minimum} degree}
\end{enumerate}
\end{tcolorbox}

\begin{tcolorbox}[colback=white, colframe=black!40, title=\textbf{Tier E --- Extension},
    fonttitle=\bfseries, arc=1.5mm, left=3mm, right=3mm, top=2mm, bottom=2mm, breakable]
\textbf{Part 1 --- The box, run backward.} The open box from the notes has volume
$V(x) = x(10-2x)(8-2x) = 4x^3 - 36x^2 + 80x$ cubic inches, with feasible domain $0 < x < 4$. A
customer wants a box holding \textbf{exactly $48$ cubic inches}.
\begin{enumerate}[label=(\alph*), leftmargin=1.8em, itemsep=6pt]
  \item Write the equation to solve, then move everything to one side and divide out the common
        factor $4$: \; \ans{$4x^3-36x^2+80x=48 \Rightarrow x^3-9x^2+20x-12=0$}
  \item List the possible rational zeros of that cubic (Lesson 4.4): \ans{$\pm1,\pm2,\pm3,\pm4,\pm6,\pm12$}
  \item Find one that works and divide it out. Depressed polynomial: \ans{$x^2-8x+12$}

        \vspace{2pt}
        Factor it: \ans{$(x-2)(x-6)$} \; All three solutions: \ans{$x=1,\,2,\,6$}
  \item Which solutions are \emph{feasible} cuts? \ans{$x=1$ and $x=2$} \; Which one must you reject, and
        why? \ans{$x=6$ --- it is outside $0<x<4$; a $6$-inch cut leaves negative side lengths}
  \item You found \emph{two} different cuts that give $48$ cubic inches. Look at the table in your
        notes and explain how one volume can happen twice. \ansline{The volume rises to about $52.5$
        near $x=1.5$ and falls again, so it passes the value $48$ once on the way up ($x=1$) and once
        on the way down ($x=2$).}
\end{enumerate}

\vspace{5pt}
\textbf{Part 2 --- Build a graph to order.} Write a polynomial in factored form whose graph
\emph{touches} the $x$-axis at $x=-1$, \emph{crosses} at $x=0$ and $x=3$, and has \textbf{both arms
pointing down}.
\begin{enumerate}[label=(\alph*), leftmargin=1.8em, itemsep=6pt]
  \item The three zeros force the factors \ans{$(x+1)$, $x$, $(x-3)$}, and the touch forces one of them to have
        exponent \ans{$2$}.
  \item So the smallest possible degree is \ans{$4$}. Is that even or odd? \ans{even}
  \item Both arms down requires an \ans{even} degree and a \ans{negative} leading coefficient.
        Does your degree from (b) cooperate? \ans{yes}
  \item Final answer: $f(x) = $ \ans{$-x(x+1)^2(x-3)$} \quad Its $y$-intercept is \ans{$(0,0)$}.
  \item Could a \emph{cubic} have done this job? \ans{no} \; Explain in one sentence.
        \ansline{A touch plus two crossings already needs degree $4$, and an odd degree can never
        point both arms the same way.}
\end{enumerate}

\vspace{5pt}
\textbf{Part 3 --- What a picture can and cannot prove.} A polynomial $p$ has real coefficients. Its
graph has exactly \textbf{two} $x$-intercepts, both simple crossings, and exactly \textbf{three}
turning points.
\begin{enumerate}[label=(\alph*), leftmargin=1.8em, itemsep=6pt]
  \item Three turning points forces the degree to be at least \ans{$4$}.
  \item The arms: with three turning points and two crossings, both arms point the \ans{same}
        direction, so the degree is \ans{even}.
  \item Take the degree to be exactly $4$. Counting with multiplicity, how many \emph{real} zeros
        does $p$ have? \ans{$2$} \; So how many \emph{imaginary} zeros? \ans{$2$} \;
        (Lesson 4.5 --- and check that the number is even.)
  \item A classmate says ``the graph proves the polynomial has degree $4$.'' Give a degree-$6$
        polynomial whose graph could look exactly the same, and say what the graph really proves.
        \ansline{$p(x)=(x+2)(x-2)(x^2+1)^2$ has the same two crossings and no new intercepts, because
        $x^2+1$ contributes only imaginary zeros. A graph proves a \emph{minimum} degree, never the
        exact degree.}
\end{enumerate}
\end{tcolorbox}

\begin{teachernote}
\textbf{Tier R} is a matching set built on one idea: three functions, the same two numbers, three
different pictures. If a group finishes item 1 quickly, make them predict the match before looking
--- the table already contains the answer, and saying so is the whole lesson. The most common Tier R
error is reading Graph I as (a) because both have a hump: send them to the arms first.

\vspace{3pt}
\textbf{Tier A} adds the two things Tier R does not require: factoring before analyzing (item 1a ---
watch for students who never pull the GCF and then cannot find $x=0$ as a zero), and solving for the
stretch factor $a$ from the $y$-intercept (item 2c). If a group writes $a=1$ in item 2 without
checking, ask them what the left arm does and let the contradiction do the teaching. Item 1(f)--(g)
is the A2.F.2a/c/d/e work: \emph{no absolute extrema} for an odd degree is the sentence to insist on.

\vspace{3pt}
\textbf{Tier E Part 1} is the payoff problem of the unit --- a modeling question that turns into a
cubic equation and is finished with the Rational Root Theorem from Lesson 4.4. Part (e) is the one
worth debriefing with the whole class: the two answers $x=1$ and $x=2$ straddle the maximum near
$x=1.5$, and both appear in the notes table as $48$. Students who noticed that in the notes will
recognize it instantly. Part 3(d) is the deepest item on the page: a graph gives a \emph{minimum}
degree, because imaginary zeros and even-multiplicity factors are invisible.

\vspace{3pt}
\textbf{Answers that look wrong but are not:} in Tier A item 1(c), the $y$-intercept is $(0,0)$ ---
the origin is both the $y$-intercept and one of the zeros. In Tier E Part 2(d), $f(0)=0$ for the same
reason.
\end{teachernote}

\end{document}
